(a) Since $y-x^2$ is irreducible, $I(Y)=(y-x^2)$, so $A(Y)=k[x,y]/(y-x^2)\cong k[x]$.

(b) $A(Z)=k[x,y]/(xy-1)\cong k[x,x^{-1}]$. The element $x$ is a nonconstant unit, whereas the only units of $k[t]$ are $k^\times$. Hence these $k$-algebras are not isomorphic.

(c) The degree-two homogeneous part of $f$ is a product of two linear forms. If they are distinct, a linear change of coordinates and scaling give $f=xy+ax+by+c$. Translating gives $(x+b)(y+a)=ab-c$, where $ab-c\ne0$ by irreducibility. Thus $A(W)\cong k[t,t^{-1}]$. If the two factors coincide, we may write $f=x^2+ax+by+c$. Here $b\ne0$, since otherwise $f$ factors over $k$. Solving for $y$ gives $A(W)\cong k[x]$. Thus the cases correspond respectively to two distinct points and one point at infinity.
